\chapter{Risque de règlement}\label{ch:m2:settlement-risk}

Le 26 juin 1974, l'autorité allemande de supervision bancaire ferma la Bankhaus Herstatt, une
banque allemande qui avait subi de lourdes pertes sur des positions de change spéculatives. Elle le fit
en pleine journée ouvrée allemande, avant l'ouverture des marchés de New York.
À ce moment-là, Herstatt avait reçu, par le système de paiement allemand, les
deutschemarks qu'elle avait achetés deux jours plus tôt~; elle n'avait pas encore payé les dollars
qu'elle avait vendus, et qui devaient être versés plus tard dans la journée à New York. Ses contreparties avaient
payé et ne seraient pas payées. Plusieurs banques furent touchées, et le principal
système américain de \omterm{def:m2:settlement-risk:netting}{compensation des paiements} en dollars ferma pendant une journée. Jusqu'alors, on pensait que le risque
d'un trader de devises se réduisait au prix~; Herstatt montra que la totalité
du principal d'une transaction pouvait être perdue entre deux paiements dans deux fuseaux horaires
(\cref{fig:m2:settlement-risk:herstatt}). Un demi-siècle plus tard,
l'enquête des banques centrales d'avril 2025 estimait qu'environ 1.4~billion
de dollars par jour continuait d'être réglé d'une manière totalement exposée à ce risque. Ce
chapitre explique le \omterm{def:m2:settlement-risk:risk}{risque de règlement}, le mécanisme \omterm{def:m2:settlement-risk:pvp}{paiement contre paiement} conçu pour
le supprimer, la \omterm{def:m2:settlement-risk:netting}{compensation des paiements} qui le réduit, et ce qui subsiste.

\section{Herstatt}

\begin{definition}[Risque de règlement, risque Herstatt]\label{def:m2:settlement-risk:risk}
Le \emph{risque de règlement}\index{risque de règlement} d'une opération de change est
le risque qu'une partie paie la devise qu'elle a vendue et ne reçoive pas la devise
qu'elle a achetée. Dans sa forme liée aux fuseaux horaires, lorsque le paiement d'une des deux jambes devient irrévocable
des heures avant l'échéance de l'autre jambe, on l'appelle \emph{risque
Herstatt}\index{risque Herstatt}.
\end{definition}

Le risque n'est pas le mouvement du marché, qui est faible sur une journée, mais le principal.
Une banque qui vend des yens contre des dollars paie ses yens le matin asiatique et reçoit
ses dollars durant l'après-midi américain~; si sa contrepartie fait défaut entre-temps, elle
a perdu les yens et n'obtient pas les dollars. L'exposition dure depuis le
moment où son propre paiement ne peut plus être rappelé jusqu'au moment où elle sait que
l'autre paiement est arrivé.

\begin{definition}[Heure de paiement unilatéral irrévocable]\label{def:m2:settlement-risk:irrevocable}
L'\emph{heure de paiement unilatéral irrévocable}\index{heure de paiement unilatéral irrévocable}
d'une devise, pour une banque donnée, est l'échéance au-delà de laquelle elle ne peut plus
annuler un paiement qu'elle a ordonné dans cette devise. L'exposition au \omterm{def:m2:settlement-risk:risk}{risque de règlement}
d'une transaction court de l'heure irrévocable de la devise vendue jusqu'à la réception
confirmée et définitive de la devise achetée, et peut être allongée par le temps
nécessaire à la vérification de la réception et par un paiement échoué.
\end{definition}

\begin{omfigure}
\centering
\begin{tikzpicture}[font=\small,x=0.42cm]
\draw[thick,-{Stealth}] (-3,0) -- (25,0) node[right,font=\footnotesize] {heures, GMT};
\foreach \h/\l in {-2/{22:00},4/{4:00},10/{10:00},16/{16:00},22/{22:00}} {\draw (\h,0.1) -- (\h,-0.1) node[below,font=\scriptsize] {\l};}
\fill[omDef!25] (-2,0.5) rectangle (7,1.0); \node[font=\footnotesize,anchor=west] at (-2,1.25) {paiements en yens (Tokyo)};
\fill[omThm!25] (6,1.8) rectangle (16,2.3); \node[font=\footnotesize,anchor=west] at (6,2.55) {paiements en euros};
\fill[omProp!25] (13,3.1) rectangle (21,3.6); \node[font=\footnotesize,anchor=west] at (13,3.85) {paiements en dollars (New York)};
\draw[omProp,very thick,{Stealth}-{Stealth}] (-2,4.7) -- (21,4.7);
\node[omProp,font=\footnotesize,anchor=south] at (9.5,4.75) {une banque qui vend des yens contre des dollars est exposée pendant environ 23 heures};
\end{tikzpicture}

\omcaption{Pourquoi les fuseaux horaires créent le risque. Les paiements dans chaque devise sont effectués
pendant les heures de leur propre système~; une banque qui vend des yens et achète des dollars ne peut pas
rappeler ses yens après l'échéance de Tokyo, et ignore qu'elle détiendra ses dollars
seulement l'après-midi à New York. Heures indicatives, GMT, le jour du
règlement.}\label{fig:m2:settlement-risk:zones}
\end{omfigure}

\begin{omfigure}
\centering
\begin{tikzpicture}[font=\small,x=0.5cm]
\draw[thick,-{Stealth}] (0,0) -- (26,0) node[right,font=\footnotesize] {temps};
\draw (2,0.1) -- (2,-0.1) node[below,font=\footnotesize,align=center] {24 juin\\transactions};
\draw (12,0.1) -- (12,-0.1) node[below,font=\footnotesize,align=center] {26 juin, midi\\Allemagne~: banque fermée};
\draw (20,0.1) -- (20,-0.1) node[below,font=\footnotesize,align=center] {Ouverture à New York~:\\dollars dus};
\draw[omDef,thick,-{Stealth}] (8,1.9) -- (11.8,0.25);
\node[omDef,font=\footnotesize,align=center] at (8,2.4) {deutschemarks reçus par Herstatt\\via le système allemand};
\draw[omProp,thick,-{Stealth}] (22,1.9) -- (20.2,0.25);
\node[omProp,font=\footnotesize,align=center] at (22.5,2.4) {dollars jamais versés\\à ses contreparties};
\end{tikzpicture}

\omcaption{La journée Herstatt, 26 juin 1974. Les contreparties avaient payé leurs deutschemarks pour
des transactions conclues deux jours plus tôt~; la banque a été fermée au milieu de la journée ouvrée allemande
avant l'ouverture de New York, et les dollars qu'elle devait n'ont pas été payés.
D'après le compte rendu de la BCE de 2007. Schéma de principe.}\label{fig:m2:settlement-risk:herstatt}
\end{omfigure}

Les banques centrales du G10 répondirent en 1996 par une stratégie en trois volets~:
les banques doivent mesurer et contrôler leurs expositions au \omterm{def:m2:settlement-risk:risk}{risque de règlement} comme
toute autre exposition de crédit~; les groupes professionnels doivent mettre au point des services qui règlent les deux jambes ensemble~;
les banques centrales doivent promouvoir ces deux démarches et adapter leurs systèmes de paiement. La réponse de la place
fut CLS, qui commença à fonctionner en septembre 2002.

\section{Paiement contre paiement}

\begin{definition}[Paiement contre paiement]\label{def:m2:settlement-risk:pvp}
Le \emph{paiement contre paiement}\index{paiement contre paiement} (PvP) est un mécanisme de règlement
dans lequel le paiement définitif d'une devise a lieu si, et seulement si, le
paiement définitif de l'autre devise a lieu.
\end{definition}

Le PvP est aux devises ce que la livraison contre paiement est aux titres
(livre~1, chapitre~5). CLS l'applique à 18 des principales devises~: elle règle
les deux jambes de chaque transaction ensemble, ou aucune. Le \omterm{def:m2:settlement-risk:risk}{risque de règlement} disparaît~;
à sa place, chaque membre doit verser à temps les devises qu'il doit, un besoin de liquidité
qui est moindre lorsque de nombreuses transactions sont réglées les unes contre les autres.

\begin{proposition}[Ce que le PvP supprime et ce qu'il laisse]\label{prop:m2:settlement-risk:pvp}
Sous PvP, les transactions d'une contrepartie défaillante ne sont pas réglées et le survivant conserve
la devise qu'il aurait versée~: sa perte est le coût du remplacement de la transaction au
nouveau cours du marché, soit le risque de remplacement, et non le principal. Le risque de remplacement est
de l'ordre de la variation du cours de change de la journée, une faible fraction du principal.
\end{proposition}

\begin{proof}
Si la jambe de la contrepartie n'est pas payée, celle du survivant ne l'est pas non plus, de sorte
qu'aucun principal ne sort de ses mains. Il doit tout de même acheter ailleurs la devise dont il avait besoin, au
cours de cet instant plutôt que du cours convenu~: la différence constitue sa perte.
\end{proof}

\section{Compensation}

\begin{definition}[Compensation des paiements]\label{def:m2:settlement-risk:netting}
Une \emph{compensation des paiements}\index{compensation des paiements} (payment netting) est le règlement, entre deux
parties et pour chaque devise et chaque date de valeur, du seul solde des montants qu'elles
se doivent mutuellement, à la place du paiement brut de chaque transaction.
\end{definition}

Une banque qui vend 400~millions d'euros contre dollars à une contrepartie et lui en rachète
300~millions d'euros à la même date de valeur n'a besoin de payer que 100~millions d'euros nets~;
son exposition passe de 400~millions d'euros à 100~millions d'euros. La BRI estime
que la compensation préalable au règlement a supprimé environ 1.3~billion de dollars de paiements par jour en
avril 2022. La compensation réduit le montant exposé au risque, mais non sa nature~: ce qui est encore payé brut peut
toujours être perdu.

\begin{example}[Journée de règlement d'une banque]\label{ex:m2:settlement-risk:book}
Le carnet d'une banque pour une date de valeur donnée compte huit transactions avec quatre contreparties,
pour 1.53~milliard de dollars au total~: avec A, elle vend 400~millions d'euros et achète 300~millions d'euros
contre dollars~; avec B, elle vend des yens pour 250~millions de dollars et en achète pour
150~millions de dollars~; avec C, elle vend des euros contre livres sterling (200~millions de livres) et des dollars contre
euros (100~millions d'euros)~; avec D, elle vend une devise de marché émergent hors CLS pour
80~millions de dollars et la rachète pour 50~millions de dollars. Réglé brut selon les
heures illustratives de la \cref{fig:m2:settlement-risk:zones}, son exposition culmine à
1.33~milliard de dollars l'après-midi, lorsque les paiements en dollars arrivent à échéance et que les paiements en euros
sont déjà réglés. Compensés par contrepartie et par devise, les paiements tombent à 430~millions
de dollars et le pic à 430~millions de dollars~; réglés par PvP dans toutes les devises sauf
la dernière, le pic n'est que de 30~millions de dollars, soit le solde des transactions de D
(\cref{fig:m2:settlement-risk:exposure}).
\end{example}

\begin{omfigure}
\begin{tikzpicture}
\begin{axis}[width=11cm,height=5cm,
  xlabel={heure de la journée de règlement (GMT)},ylabel={exposition (M USD)},
  tick label style={font=\small},label style={font=\small},
  xmin=-4,xmax=24,ymin=0,ymax=1450,grid=major,grid style={gray!25},
  legend columns=3,legend style={font=\footnotesize,at={(0.5,-0.3)},anchor=north,draw=none,fill=none,/tikz/every even column/.append style={column sep=8pt}},
  table/col sep=comma]
\addplot[omProp,very thick,const plot] table[x=t,y=gross]{figdata/markets-2/20-settlement-risk/exposure.csv};
\addplot[omThm,very thick,dashed,const plot] table[x=t,y=net]{figdata/markets-2/20-settlement-risk/exposure.csv};
\addplot[omDef,very thick,dotted,const plot] table[x=t,y=pvp]{figdata/markets-2/20-settlement-risk/exposure.csv};
\legend{brut,compensé,compensé et PvP}
\end{axis}
\end{tikzpicture}

\omcaption{L'exposition au règlement du carnet de
l'\cref{ex:m2:settlement-risk:book} au fil de la journée de règlement. Le règlement brut
expose la banque à l'essentiel de ses achats de la journée à la fois, l'après-midi~; la compensation
réduit l'exposition de deux tiers~; le PvP ne laisse subsister que la devise qu'il ne couvre pas.
Illustratif~; données~: tutoriel du chapitre.}\label{fig:m2:settlement-risk:exposure}
\end{omfigure}

\section{Ce qui reste non protégé}

Le \omterm{def:m2:settlement-risk:risk}{risque de règlement} n'a pas disparu. En avril 2022, selon l'estimation de la BRI, 2.2~billions de dollars du
volume de change livrable de la journée, soit près d'un tiers, ont encore été réglés à risque~; dans les
principales places de négociation, 20--40~\% du volume, dans certaines places plus petites plus de
trois quarts. L'enquête de 2025 a mesuré le phénomène autrement, selon la méthode de règlement~: un peu plus
d'un tiers du règlement de la journée, soit 5.2~billions de dollars, est passé par le PvP~; 7.6~billions de dollars ont recouru à des
méthodes qui atténuent le risque sans le supprimer, comme la compensation, le règlement au sein d'un groupe ou les comptes bancaires
assortis de contrôles de calendrier~; et 1.4~billion de dollars, soit 10~\%, a été réglé brut entre les deux
parties. Les principales raisons invoquées étaient que la contrepartie n'avait pas accès au PvP ou que la devise
n'était pas prise en charge (\cref{fig:m2:settlement-risk:methods}).
supported (\cref{fig:m2:settlement-risk:methods}).

\begin{omfigure}
\begin{tikzpicture}
\begin{axis}[width=11cm,height=4.6cm,xbar,bar width=12pt,
  xlabel={billions de dollars par jour, avril 2025},
  tick label style={font=\small},label style={font=\small},
  xmin=0,xmax=9.2,ymin=-0.6,ymax=2.6,grid=major,grid style={gray!25},
  ytick={0,1,2},yticklabels={bilatéral brut,{compensation, intragroupe, contrôles de calendrier},paiement contre paiement},
  yticklabel style={font=\footnotesize,text width=4.2cm,align=right},
  nodes near coords,every node near coord/.append style={font=\scriptsize},
  table/col sep=comma]
\addplot[fill=omThm!60,draw=omThm] coordinates {(1.4,0) (7.6,1) (5.2,2)};
\end{axis}
\end{tikzpicture}

\omcaption{Comment le règlement des changes mondial a été effectué en avril 2025, selon la méthode~: un peu plus d'un
tiers par \omterm{def:m2:settlement-risk:pvp}{paiement contre paiement}, plus de la moitié par des méthodes qui réduisent mais ne suppriment pas
le \omterm{def:m2:settlement-risk:risk}{risque de règlement}, et un dixième, 1.4~billion de dollars par jour, brut entre les deux
parties. Données~: enquête triennale de la BRI, telle que publiée dans la \emph{BIS Quarterly
Review}, juin 2026.}\label{fig:m2:settlement-risk:methods}
\end{omfigure}

Mesurer sa propre exposition est en soi une discipline. L'exposition d'une banque à une transaction
ne prend pas fin lorsque la devise achetée est due, mais lorsqu'elle sait qu'elle est arrivée~;
la BCE décrivait, en 2007, la période irrévocable qui se termine avec la réception finale, une
période incertaine tant que la banque ne l'a pas vérifiée, et une
période d'échec si le paiement n'est pas arrivé. Une banque qui rapproche ses paiements entrants seulement le
lendemain matin porte l'ensemble de ses achats pendant la nuit dans la période incertaine, et un
paiement défaillant prolonge l'exposition jusqu'à sa résolution.

Les pertes ne relèvent pas que de l'histoire. L'article de la BRI de 2022 rappelait la perte de la KfW de
300~millions d'euros lors de la faillite de Lehman Brothers en 2008, et la perte de Barclays de 130~millions
de dollars face à une petite société de change en mars 2020. Le schéma est celui de Herstatt~: un
paiement effectué dans l'hypothèse que l'autre suivrait.

\begin{dated}{2026-09}{Le règlement aujourd'hui}\label{dat:m2:settlement-risk:today}
Enquête triennale de la BRI, avril 2025~: 5.2~billions de dollars par jour réglés par PvP, 7.6~billions de dollars par des
méthodes qui atténuent le risque, 1.4~billion de dollars brut bilatéral (10~\%).
CLS règle 18 devises, en moyenne plus de 8~billions de dollars par jour, pour plus de
75 membres de règlement.
\end{dated}

\section{Tutoriel~: la chronologie de l'exposition au règlement}\label{tut:m2:settlement-risk:timeline}

\begin{tutorial}
\textbf{Objectif.} Calculer l'exposition au règlement d'un carnet au fil de la journée, brute, compensée
et avec PvP, à partir de l'échéance d'annulation et du délai de confirmation de chaque devise.
\textbf{État final~:} \cref{fig:m2:settlement-risk:exposure},
\cref{ex:m2:settlement-risk:book} et les chiffres du problème du week-end.
\begin{tutsteps}
\item \textbf{La fenêtre d'exposition} d'une transaction et l'exposition à un instant donné.
\omcode{code/firm/settlerisk/firm_settlerisk.py}{20}{28}{Fenêtre d'exposition d'une transaction et exposition d'un carnet à un instant donné.}\label{lst:m2:settlement-risk:window}
\item \textbf{Compensation et profil de la journée}, les devises PvP étant retirées.
\omcode{code/firm/settlerisk/firm_settlerisk.py}{31}{61}{Compensation des paiements par contrepartie, et profil d'exposition de la journée.}\label{lst:m2:settlement-risk:profile}
\item \textbf{Exécuter} \texttt{settle\_demo.summary()} et \texttt{fig\_settle.py}.
\end{tutsteps}
\textbf{Pour aller plus loin.} Ajouter une période incertaine, les heures pendant lesquelles la banque
confirme une réception, à chaque devise, ainsi qu'un paiement défaillant qui prolonge l'exposition d'une transaction
jusqu'au lendemain~; observer comment le pic se déplace.
\end{tutorial}

\section{Construction~: le calculateur d'exposition au risque de règlement}\label{bld:m2:settlement-risk:settlerisk}

\begin{build}
\bfield{Rôle} La petite société règle ses propres transactions de change et celles de ses
clients~: elle doit savoir, heure par heure, combien elle pourrait perdre si une
contrepartie faisait défaut, fixer des limites en conséquence, et voir ce que la \omterm{def:m2:settlement-risk:netting}{compensation des paiements} et le PvP permettraient d'économiser.
\bfield{Interface} \texttt{Trade(counterparty, sell, buy, value\_usd)}~;
\texttt{window(trade, times)}~; \texttt{exposure\_at(t, trades, times)}~;
\texttt{net\_by\_counterparty(trades)}~; \texttt{profile(trades, times, pvp)}~;
\texttt{peak(profile)}.
\bfield{Règles} Heures par devise, sous la forme (échéance d'annulation, confirmation), en
heures GMT du jour de règlement~; montants en dollars~; \omterm{def:m2:settlement-risk:netting}{compensation des paiements} par
contrepartie et par devise~; les transactions entre deux devises PvP ne génèrent aucune exposition au règlement.
\bfield{Tests d'acceptation} \texttt{code/firm/settlerisk/tests/}~: la fenêtre de vente de yens
contre dollars et la fenêtre inverse, vide~; exposition sommée et nulle après la dernière
confirmation~; le PvP la supprime~; la compensation réduit les paiements et le pic.
\bfield{Pour aller plus loin} Limites et alertes par contrepartie~; périodes incertaines et défaillantes~;
liquidité PvP~: les apports de liquidité nécessaires à un membre, par devise et par heure.
\end{build}

\begin{omsources}
\item European Central Bank, \emph{Financial Stability Review}, décembre 2007,
Encadré 19, ``More than thirty years after the Herstatt case''.
\item M.~Glowka et T.~Nilsson, ``FX settlement risk: an unsettled issue'',
\emph{BIS Quarterly Review}, décembre 2022.
\item M.~Drehmann, P.~McGuire, T.~Shirakami, M.~Conway et N.~Lovell, ``Uncovering FX settlement risk: new measures from the 2025 BIS Triennial Survey'',
\emph{BIS Quarterly Review}, juin 2026.
\item CLS Group, informations sur la société.
\end{omsources}

\section{Exercices}

\begin{exercise}[$\star$]\label{exo:m2:settlement-risk:1}
Avec les heures illustratives de la \cref{fig:m2:settlement-risk:zones}, calculer pendant
combien de temps une banque est exposée lorsqu'elle vend des yens contre des dollars et lorsqu'elle vend des dollars contre des yens.
\end{exercise}

\begin{exercise}[$\star$]\label{exo:m2:settlement-risk:2}
Calculer la perte d'une banque si sa contrepartie fait défaut pendant la fenêtre
d'exposition, dans le cadre d'un règlement brut, puis dans celui du PvP.
\end{exercise}

\begin{exercise}[$\star$]\label{exo:m2:settlement-risk:3}
Deux banques ont dix transactions en EURUSD pour une même date de valeur~: la première vend
EUR~1 milliard au total et achète EUR~900 millions. Calculer le montant versé au titre de la \omterm{def:m2:settlement-risk:netting}{compensation des paiements}.
\end{exercise}

\begin{exercise}[$\star\star$]\label{exo:m2:settlement-risk:4}
Dans l'\cref{ex:m2:settlement-risk:book}, calculer les paiements nets avec chaque contrepartie.
\end{exercise}

\begin{exercise}[$\star\star$]\label{exo:m2:settlement-risk:5}
Expliquer pourquoi le PvP crée un besoin de liquidité, et comment le règlement conjoint
de nombreuses transactions le réduit.
\end{exercise}

\begin{exercise}[$\star\star$]\label{exo:m2:settlement-risk:6}
Expliquer pourquoi le \omterm{def:m2:settlement-risk:risk}{risque de règlement} est plus élevé pour les devises des pays émergents
et auprès des plus petites contreparties.
\end{exercise}

\begin{exercise}[$\star\star\star$]\label{exo:m2:settlement-risk:7}
\emph{Programmation.} Avec \texttt{profile}, calculer l'exposition du carnet à 05:00,
08:00, 14:00 et 17:00, brute et après compensation.
\end{exercise}

\begin{exercise}[$\star\star\star$]\label{exo:m2:settlement-risk:8}
\emph{Repérer l'erreur.} «~Notre exposition de change à une contrepartie est la valeur de marché de nos
contrats à terme de change ouverts avec elle, soit quelques millions~; le \omterm{def:m2:settlement-risk:risk}{risque de règlement} est négligeable.» Corriger ce raisonnement.
\end{exercise}

\section{Problème~: Le trillion non réglé}

\begin{problem}[{Problème du week-end --- la journée de règlement d'une banque}]\label{pb:m2:settlement-risk:1}
Une banque dispose du carnet de l'\cref{ex:m2:settlement-risk:book} pour la date de valeur du lendemain,
avec les heures illustratives de paiement de la \cref{fig:m2:settlement-risk:zones} (yen
de 22:00 la veille à 07:00, euro et livre sterling 06:00--16:00 et
07:00--16:00, dollar 13:00--21:00, la devise émergente 04:00--12:00).

\medskip\noindent\textbf{Partie I --- Brut.}
\begin{enumerate}
\item Donner les achats bruts de la journée et le pic d'exposition.
\item Quand se situe le pic, et pourquoi à ce moment-là~?
\item Quelle transaction isolée est exposée le plus longtemps, et pendant combien de temps~?
\item Quelle est l'exposition à 05:00~?
\item Si la contrepartie A faisait défaut à 14:00, quelle perte la banque subirait-elle~?
\end{enumerate}

\medskip\noindent\textbf{Partie II --- Compensation.}
\begin{enumerate}[resume]
\item Donner les paiements nets avec chaque contrepartie.
\item Donner le pic après compensation.
\item Quelle perte la banque subirait-elle si A faisait défaut à 14:00 après compensation~?
\item Que faut-il sur le plan juridique pour que la compensation tienne en cas de défaillance~?
\item Que reste-t-il avec D après compensation, et pourquoi~?
\end{enumerate}

\medskip\noindent\textbf{Partie III --- PvP.}
\begin{enumerate}[resume]
\item Donner le pic avec PvP pour toutes les devises sauf la devise émergente.
\item Quel risque la banque court-elle désormais avec A, B et C~?
\item Quel est le coût du PvP pour la banque~?
\item Comment pourrait-on protéger les 30~millions de dollars restants~?
\item Pourquoi de nombreuses banques règlent-elles encore certaines transactions en brut~?
\end{enumerate}

\medskip\noindent\textbf{Partie IV --- Jugement.}
\begin{enumerate}[resume]
\item Comment une banque doit-elle fixer une limite de règlement pour une contrepartie~?
\item Qu'est-ce que l'affaire Herstatt a changé dans la façon dont les banques perçoivent le risque de change~?
\item Que signifie le chiffre de 10~\% de l'enquête 2025 pour l'ensemble du marché~?
\item Énoncer le \emph{résultat nommé}~: le pic d'exposition du carnet en brut, après compensation et
avec PvP.
\item En une phrase~: qu'est-ce que le \omterm{def:m2:settlement-risk:risk}{risque de règlement}, et qu'est-ce qui le supprime~?
\end{enumerate}
\end{problem}

\section{Questions d'entretien}

\begin{interviewq}[$\star$ \iqroles{banque, développeur}]\label{iq:m2:settlement-risk:1}
Qu'est-ce que le \omterm{def:m2:settlement-risk:risk}{risque Herstatt}, et comment CLS le supprime-t-il~?
\end{interviewq}

\begin{interviewq}[$\star$ \iqroles{banque}]\label{iq:m2:settlement-risk:2}
Quelle est la différence entre le \omterm{def:m2:settlement-risk:risk}{risque de règlement} et le risque de remplacement~?
\end{interviewq}

\begin{interviewq}[$\star\star$ \iqroles{développeur}]\label{iq:m2:settlement-risk:3}
Comment calculeriez-vous l'exposition intrajournalière au règlement d'une banque vis-à-vis de chaque contrepartie~?
\end{interviewq}

\begin{interviewq}[$\star\star$ \iqroles{banque, trader}]\label{iq:m2:settlement-risk:4}
Pourquoi tout le change n'est-il pas réglé par le PvP~?
\end{interviewq}

\begin{interviewq}[$\star\star$ \iqroles{chercheur, banque}]\label{iq:m2:settlement-risk:5}
Dans quelle mesure la \omterm{def:m2:settlement-risk:netting}{compensation des paiements} peut-elle réduire l'exposition au règlement, et qu'est-ce qui la limite~?
\end{interviewq}

\begin{interviewq}[$\star\star\star$ \iqroles{développeur, banque}]\label{iq:m2:settlement-risk:6}
Concevoir un système qui bloque les nouvelles transactions avec une contrepartie lorsque son
exposition projetée au règlement à une date de valeur dépasserait sa limite.
\end{interviewq}
